\section{Results}\label{sec:results}
\subsection{From baseline to competence}
The base models solve none of the counted slots of the 100 exams, with four reruns of every failed task, at 0.6B, 4B and 8B
(\cref{fig:exam}); \uncountedFourBase{} base slots at 4B and \uncountedEightBase{} at 8B never counted. After training on the round-one set the 4B adapter solves \examFourFullOne{} tasks on the first
attempt and \examFourFullFive{} within five attempts (95\% interval \examFourFullFiveLo--\examFourFullFiveHi{}\%), and
the 8B adapter \examEightFullOne{} and \examEightFullFive{} (\examEightFullFiveLo--\examEightFullFiveHi{}\%). Over the
tasks both arms counted, the paired sign test finds \winsFourFull{} wins and \lossesFourFull{} losses at 4B and
\winsEightFull{} wins and \lossesEightFull{} losses at 8B (one-sided $p$ \pFourFull{} and \pEightFull{}). The 0.6B
adapter solves \examZeroFullFive{} task within five attempts, so we treat that size as not reaching the exams with this tool
and data, although on the practice tasks it goes from \trainZeroBasePassed{} to \trainZeroFullPassed{} passes in
\trainZeroBaseTrials{} attempts. On the practice tasks the 4B base passes \trainFourBasePassed{} of
\trainFourBaseTrials{} attempts and its adapter \trainFourFullPassed{}; the 8B base \trainEightBasePassed{} and its
adapter \trainEightFullPassed{} (\cref{fig:train}). \Cref{tab:stages} lists every stage and \cref{fig:loss} the held-out loss of training.

A pre-restart measurement of an adapter trained the same way on the same records solved \examFourFirstFive{} (4B) and
\examEightFirstFive{} (8B) within five attempts against \examFourFullFive{} and \examEightFullFive{}, and \examFourFirstOne{} and
\examEightFirstOne{} on the first attempt against \examFourFullOne{} and \examEightFullOne{}: the first attempt of 100 tasks moved by four tasks at 4B and eleven at 8B between
measurements, and we therefore treat solved-within-five as the measurement and pass@1 as a noisy companion.

\begin{figure}[t]
\centering
\begin{tikzpicture}
\begin{groupplot}[group style={group size=2 by 1, horizontal sep=14mm}, paper, width=0.5\linewidth, height=4.6cm,
  symbolic x coords={0.6B,4B,8B}, xtick=data, ymin=0, ymax=40, enlarge x limits=0.25,
  ylabel={tasks of 100}, legend style={at={(0.03,0.97)}, anchor=north west}]
\nextgroupplot[ybar, bar width=11pt, title={pass@1}]
\addplot[fill=okGrey!60, draw=black!70] coordinates {(0.6B,\examZeroBaseOne) (4B,\examFourBaseOne) (8B,\examEightBaseOne)};
\addplot[fill=okBlue!80, draw=black!70] coordinates {(0.6B,\examZeroFullOne) (4B,\examFourFullOne) (8B,\examEightFullOne)};
\legend{base, adapter}
\nextgroupplot[ybar, bar width=11pt, title={solved within five attempts}]
\addplot[fill=okGrey!60, draw=black!70] coordinates {(0.6B,\examZeroBaseFive) (4B,\examFourBaseFive) (8B,\examEightBaseFive)};
\addplot[fill=okBlue!80, draw=black!70] coordinates {(0.6B,\examZeroFullFive) (4B,\examFourFullFive) (8B,\examEightFullFive)};
\end{groupplot}
\end{tikzpicture}
\caption{Held-out exams, 100 tasks. The base models solve none, even with four reruns of every failed task. The adapters
  were trained on the round-one set; pass@1 is the first attempt of each task.}\label{fig:exam}
\end{figure}

\begin{figure}[t]
\centering
\begin{tikzpicture}
\begin{groupplot}[group style={group size=2 by 1, horizontal sep=14mm}, paper, width=0.5\linewidth, height=4.4cm,
  xlabel={attempts allowed}, xtick={1,2,3,4,5}, ymin=0, ymax=34, ylabel={tasks solved}, legend style={at={(0.97,0.03)}, anchor=south east, font=\scriptsize}]
\nextgroupplot[title={4B}]
\addplot[okGrey, thick, mark=square*, mark size=1.6pt] table[x=k,y=base] {withink-four.dat};
\addplot[okBlue, thick, mark=*, mark size=1.6pt] table[x=k,y=full] {withink-four.dat};
\addplot[okOrange, thick, dashed, mark=triangle*, mark size=2pt] table[x=k,y=first] {withink-four.dat};
\legend{base, adapter, pre-restart measurement}
\nextgroupplot[title={8B}]
\addplot[okGrey, thick, mark=square*, mark size=1.6pt] table[x=k,y=base] {withink-eight.dat};
\addplot[okBlue, thick, mark=*, mark size=1.6pt] table[x=k,y=full] {withink-eight.dat};
\addplot[okOrange, thick, dashed, mark=triangle*, mark size=2pt] table[x=k,y=first] {withink-eight.dat};
\end{groupplot}
\end{tikzpicture}
\caption{Exam tasks solved as the number of allowed attempts grows (the first attempt, then reruns of the failed tasks). The dashed curve is a pre-restart measurement of an adapter trained the same way on the same records: it differs by
  four tasks at 4B and two at 8B within five attempts, and by four and eleven on pass@1.}\label{fig:withink}
\end{figure}

\begin{figure}[t]
\centering
\begin{tikzpicture}
\begin{axis}[paper, ybar, bar width=9pt, width=0.62\linewidth, height=4.4cm, symbolic x coords={0.6B,4B,8B}, xtick=data, ymin=0, ymax=100,
  enlarge x limits=0.25, ylabel={passed attempts (\%)}, legend style={at={(0.03,0.97)}, anchor=north west}]
\addplot[fill=okGrey!60, draw=black!70] coordinates {(0.6B,\trainZeroBasePct) (4B,\trainFourBasePct) (8B,\trainEightBasePct)};
\addplot[fill=okBlue!80, draw=black!70] coordinates {(0.6B,\trainZeroFullPct) (4B,\trainFourFullPct) (8B,\trainEightFullPct)};
\legend{base, adapter}
\end{axis}
\end{tikzpicture}
\caption{Practice and composite tasks, \trainFourBaseTrials{} attempts per model (228 tasks, four attempts each). These are the
  tasks the adapters train on: the figure shows fit, not generalisation.}\label{fig:train}
\end{figure}

\begin{figure}[t]
\centering
\begin{tikzpicture}
\begin{axis}[paper, width=0.62\linewidth, height=4.6cm, ymode=log, xlabel={optimiser step}, ylabel={held-out loss},
  legend style={at={(0.97,0.97)}, anchor=north east}]
\addplot[okVermillion, thick, mark=*, mark size=1.3pt] table[x=step,y=loss] {curve-zero-full.dat};
\addplot[okBlue, thick, mark=square*, mark size=1.3pt] table[x=step,y=loss] {curve-four-full.dat};
\addplot[okGreen, thick, mark=triangle*, mark size=1.6pt] table[x=step,y=loss] {curve-eight-full.dat};
\addplot[okBlue, thick, dashed] table[x=step,y=loss] {curve-four-rtwo.dat};
\addplot[okGreen, thick, dashed] table[x=step,y=loss] {curve-eight-rtwo.dat};
\legend{0.6B, 4B, 8B, 4B integrated, 8B integrated}
\end{axis}
\end{tikzpicture}
\caption{Held-out loss while training, step 0 the base model. Solid: the round-one set; dashed: the round-two set. The
  held-out records are whole tasks the training set does not contain.}\label{fig:loss}
\end{figure}

\begin{table}[!ht]
\centering\small\setlength{\tabcolsep}{3.5pt}
\input{stages-table}
\caption{Exam results by stage. Each row is 100 tasks, one attempt and up to four reruns of the failed; loop-stopped is the
  share of attempts the watchdog ended and repeat after error the share of calls after a Vim error that repeated the failed call.
  The two rows marked $\dagger$, and the two marked $\ddagger$, were measured together in one run; the full-set adapter
  appears again because each comparison re-measures it.}\label{tab:stages}
\end{table}

\subsection{What the data does}
\paragraph{The teacher solved \recoverySolvedPct{}\% of the continuations it attempted.} The adapters were sent back at the practice tasks
(\recoveryAttempts{} attempts over three sizes, \recoveryPassPct{}\% passed). The teacher read each task whose attempt
failed, continued \recoveryContinued{} of those attempts from the step where they went wrong, and \recoverySolved{} of the
continuations verified; each is training data that shows the student its own mistake and the call that fixes it. No arm of this
study removes the continuations, so their effect on the student is not isolated. Of the \integrationCandidates{} candidate records the round
produced, \integrationAdmitted{} were new solutions of their task; \integrationDuplicate{} repeated a tool-call sequence the
task already had and \integrationOvercap{} exceeded the per-task cap of sixteen (\cref{tab:data}).

\paragraph{The round-two set gave the 8B a larger count that the test does not confirm, and the 4B none.} Trained on
the round-two set (\integrationTrainBefore{} to \integrationTrainAfter{} training records), the 8B adapter solves \examEightRtwoFive{} tasks within five attempts against \examEightRtworefFive{}
for the round-one adapter measured in the same run (\winsEightRtwo{} wins, \lossesEightRtwo{} losses, one-sided $p$
\pEightRtwo{}); the 4B adapter solves \examFourRtwoFive{} against \examFourRtworefFive{} (\winsFourRtwo{} wins,
\lossesFourRtwo{} losses). The 8B count is higher by nine tasks but the one-sided test does not reach significance; the 4B shows no gain.
Both are single measurements, and pass@1 at the 8B moved the other way (\examEightRtwoOne{} against
\examEightRtworefOne{}).

\subsection{How the models fail}
After a Vim error the base models most often repeat the failed call: \behaviourFourBaseExamRepeatPct{}\%
of the 4B base's next calls and \behaviourEightBaseExamRepeatPct{}\% of the 8B base's repeat the call that just failed
(\behaviourZeroBaseExamRepeatPct{}\% at 0.6B). Training cuts this to \behaviourFourFullExamRepeatPct{}\% at 4B and
\behaviourEightFullExamRepeatPct{}\% at 8B, the 0.6B stays at \behaviourZeroFullExamRepeatPct{}\%. A run that repeats
itself is ended by the loop watchdog after six identical calls. That share did not fall with training:
\behaviourFourFullExamLoopPct{}\% of the 4B adapter's exam attempts and \behaviourEightFullExamLoopPct{}\% of the 8B adapter's end
that way, against \behaviourFourBaseExamLoopPct{}\% and \behaviourEightBaseExamLoopPct{}\% for the bases, although the share of
repeated failed calls did. Runs that never leave a loop are therefore a large share of the failures that remain; we did not classify the rest.

\subsection{A loop penalty in the policy gradient}
Starting from the round-one adapter, a clipped policy gradient on its own graded attempts (eight attempts of every practice
task at temperature 1.0), with the penalties of \cref{sec:method} (a full pass's reward for a loop-stopped attempt, half a
pass for repeating failed calls), moved the policy (median per-step KL to its starting adapter \optimiserFourKl{} at 4B and
\optimiserEightKl{} at 8B, median clip fraction \optimiserFourClip{}\% and \optimiserEightClip{}\%). Against the adapter it started from, the share of exam attempts
ended by the loop watchdog is \behaviourFourPgExamLoopPct{}\% against \behaviourFourPgrefExamLoopPct{}\% at 4B and
\behaviourEightPgExamLoopPct{}\% against \behaviourEightPgrefExamLoopPct{}\% at 8B; the share of repeated failed calls, which the penalty
targets directly, rose at both sizes: \behaviourFourPgExamRepeatPct{}\% against \behaviourFourPgrefExamRepeatPct{}\% and
\behaviourEightPgExamRepeatPct{}\% against \behaviourEightPgrefExamRepeatPct{}\%. The 4B solves \examFourPgFive{} tasks
within five attempts against \examFourPgrefFive{}. The 8B solves \examEightPgFive{} against
\examEightPgrefFive{} (\winsEightPg{} wins, \lossesEightPg{} losses, $p$ \pEightPg{}), a gain that comes with a loop share only four points lower, so it is not evidence that the penalty works.
A penalty twice as strong (\optimiserFourStrongLoop{}\% and \optimiserFourStrongRepeat{}\% of a pass's reward instead of
\optimiserFourLoop{}\% and \optimiserFourRepeat{}\%), trained for twice as many steps (median KL \optimiserFourStrongKl{}, clip fraction
\optimiserFourStrongClip{}\%), raised the loop-stopped share of the 4B to \behaviourFourPgdExamLoopPct{}\% against
\behaviourFourPgdrefExamLoopPct{}\% for the adapter it started from and solved \examFourPgdFive{} tasks within five attempts
against \examFourPgdrefFive{} (\winsFourPgd{} wins, \lossesFourPgd{} losses); its first-attempt count fell from
\examFourPgdrefOne{} to \examFourPgdOne{}.
On the practice tasks looping is rare (\behaviourFourFullTrainLoopPct{}\% of the 4B adapter's attempts and
\behaviourEightFullTrainLoopPct{}\% of the 8B's, against about 30\% on the exams), which leaves the penalty few
examples of the behaviour it targets.
